\documentclass[11pt]{article}
\usepackage[a4paper,margin=2.5cm]{geometry}
\usepackage{amsmath,amssymb,amsthm}
\usepackage[colorlinks=true,linkcolor=blue,urlcolor=blue]{hyperref}

\newcommand{\F}{\mathbb{F}}
\newcommand{\Z}{\mathbb{Z}}
\theoremstyle{plain}
\newtheorem{conj}{Conjecture}
\newtheorem{prop}{Proposition}
\theoremstyle{remark}
\newtheorem{remark}{Remark}

\title{Disproof of TLMC Conjecture 00000001040\\
\large A counterexample at the prime $q = 5$}
\author{TLMC submission}
\date{October 2, 2026}

\begin{document}
\maketitle

\begin{conj}[00000001040]
The complete-mapping polynomials (polynomials $f$ with $f(x)+x$ a permutation)
at $q \equiv 2 \pmod 3$ form the complete equivalence class of $x^3$ type
(a cubic classification of complete mappings).
\end{conj}

\begin{prop}[Counterexample]\label{prop:main}
Let $q = 5$, which is prime and satisfies $5 \equiv 2 \pmod 3$. Then:
\begin{enumerate}
\item $f(x) = 2x$ is a complete mapping of $\F_5$: both $2x$ and
$f(x)+x = 3x$ are permutations, with
$2x \mapsto (0,2,4,1,3)$ and $3x \mapsto (0,3,1,4,2)$;
\item $x^3 + x$ is not a permutation of $\F_5$: it sends $x = 0, 2, 3$ all
to $0$ (values $(x^3+x)_{x=0}^{4} = (0,2,0,0,3)$);
\item more strongly, \emph{no} member of the affine $x^3$-family
$A(x+B)^3 + C$ with $A \in \F_5^\times$, $B, C \in \F_5$ is a complete
mapping;
\item nevertheless exactly $15$ linear complete mappings $ax+b$ exist, with
slopes $a \in \{1,2,3\}$ and $b \in \F_5$, and a brute-force census shows
$\Z_5$ has exactly $15$ complete mappings in total.
\end{enumerate}
Hence at $q = 5$ the $x^3$-type class contains \emph{no} complete mapping
while $15$ complete mappings exist: the classification of
Conjecture~00000001040 fails. \qed
\end{prop}

\begin{proof}
The arithmetic claims are finite computations over $\F_5$; we verify them by
hand and mechanically.

(1) $(2x)_{x=0}^{4} = (0,2,4,1,3)$ and $(3x)_{x=0}^{4} = (0,3,1,4,2)$ are
both permutations of $\{0,1,2,3,4\}$, so $f=2x$ satisfies the complete
mapping condition $f + x$ = $3x$ permutation.

(2) $(x^3+x)_{x=0}^{4} = (0,\,2,\,10,\,30,\,68) \equiv (0,2,0,0,3) \pmod 5$:
the values at $x = 0, 2, 3$ collide at $0$, so $x^3$ is not a complete
mapping at $q = 5$.

(3) For $f = A(x+B)^3 + C$ we have $f + x = A y^3 + y + (C - B)$ with
$y = x + B$; since adding a constant and precomposing with the bijection
$x \mapsto x+B$ preserve permutation-ness, $f$ is a complete mapping iff
$A y^3 + y$ is a permutation. Direct evaluation for $A \in \F_5^\times$:
\[
\begin{array}{ll}
A=1: & (0,2,0,0,3) \text{ --- } 0,2,3 \text{ collide;}\\
A=2: & (0,3,3,2,2) \text{ --- } 1,2 \text{ and } 3,4 \text{ collide;}\\
A=3: & (0,4,1,4,1) \text{ --- } 1,3 \text{ and } 2,4 \text{ collide;}\\
A=4: & (0,0,4,1,0) \text{ --- } 0,1 \text{ and } 4 \text{ repeat } 0.
\end{array}
\]
None is a permutation, so the entire affine $x^3$-class is free of complete
mappings at $q = 5$.

(4) $ax + b$ is a permutation iff $a \neq 0$, and $(a+1)x+b$ is a
permutation iff $a + 1 \neq 0 \pmod 5$; hence $a \in \{1,2,3\}$,
$b \in \F_5$: exactly $3 \cdot 5 = 15$ linear complete mappings. A
brute-force census over all $5! = 120$ permutations of $\F_5$ (every
function $\F_5 \to \F_5$ is a polynomial of degree $< 5$) confirms that the
total number of complete mappings of $\Z_5$ is exactly $15$, i.e.\ precisely
the linear ones.

Since complete mappings exist at $q = 5$ but the $x^3$-type equivalence
class contains none (and, \emph{a fortiori}, the linear complete mappings
are not cubic-type under any degree-preserving equivalence such as affine
or EA-equivalence), the classification claimed by the conjecture is false.
\end{proof}

\begin{remark}[Boundary]
\begin{itemize}
\item $q = 5$ is the smallest non-degenerate prime power with
$q \equiv 2 \pmod 3$ ($q = 2$ is degenerate: $\F_2$ admits no complete
mapping at all, so the conjecture holds vacuously there).
\item The same obstruction persists at every prime $q \equiv 2 \pmod 3$,
$q \geq 5$: the linear map $f = 2x$ has $f + x = 3x$ a permutation
(since $3 \nmid q$), and being of degree $1$ it cannot lie in the
degree-$3$ $x^3$ equivalence class; meanwhile $x^3$ itself is a complete
mapping iff $x^3 + x$ is a permutation (e.g.\ it is one at $q = 11$, where
$-1$ is a quadratic non-residue), which does not rescue a classification
that must also include the linear complete mappings.
\end{itemize}
\end{remark}

\section*{Mechanical verification}

The bundled artifacts reproduce every number above without third-party
dependencies:
\begin{itemize}
\item \texttt{reproduce.py} --- independent Python recomputation (exits $0$
iff all checks pass; prints the five value-tuples listed in the proof);
\item \texttt{lean4/} --- a Mathlib-free Lean~4 (v4.33.1) development whose
main theorem \texttt{TLMC1040.disproof\_00000001040} asserts the
conjunction of the numerical facts (1)--(4) together with the explicit
$25$-case classification of the linear candidates
($a \in \{0,1,2,3,4\} \times b \in \F_5$), all discharged by
\texttt{decide}. \texttt{lake env lean Check.lean} audits every theorem
with \texttt{\#print axioms}: all report ``does not depend on any
axioms'' (no \texttt{sorry}, no \texttt{native\_decide}).
\end{itemize}

\end{document}
